\documentclass[10pt,a4paper]{amsart}
\usepackage{amsmath,amssymb,amsthm,graphicx}
\usepackage{verbatim}
\newenvironment{reference}{\comment}{\endcomment}
%\newenvironment{reference}{}{}
\newenvironment{slogan}{\comment}{\endcomment}
\newenvironment{history}{\comment}{\endcomment}
% For commutative diagrams we use Xy-pic
\usepackage[all]{xy}

% We use 2cell for 2-commutative diagrams.
\xyoption{2cell}
% Theorem environments.
%
\theoremstyle{plain}
\newtheorem{theorem}[subsection]{Theorem}
\newtheorem{proposition}[subsection]{Proposition}
\newtheorem{lemma}[subsection]{Lemma}

\theoremstyle{definition}
\newtheorem{definition}[subsection]{Definition}
\newtheorem{example}[subsection]{Example}
\newtheorem{exercise}[subsection]{Exercise}
\newtheorem{situation}[subsection]{Situation}

\theoremstyle{remark}
\newtheorem{remark}[subsection]{Remark}
\newtheorem{remarks}[subsection]{Remarks}

\numberwithin{equation}{subsection}

% Macros
%
\def\lim{\mathop{\mathrm{lim}}\nolimits}
\def\colim{\mathop{\mathrm{colim}}\nolimits}
\def\Spec{\mathop{\mathrm{Spec}}}
\def\Hom{\mathop{\mathrm{Hom}}\nolimits}
\def\Ext{\mathop{\mathrm{Ext}}\nolimits}
\def\SheafHom{\mathop{\mathcal{H}\!\mathit{om}}\nolimits}
\def\SheafExt{\mathop{\mathcal{E}\!\mathit{xt}}\nolimits}
\def\Sch{\mathit{Sch}}
\def\Mor{\mathop{\mathrm{Mor}}\nolimits}
\def\Ob{\mathop{\mathrm{Ob}}\nolimits}
\def\Sh{\mathop{\mathit{Sh}}\nolimits}
\def\NL{\mathop{N\!L}\nolimits}
\def\CH{\mathop{\mathrm{CH}}\nolimits}
\def\proetale{{pro\text{-}\acute{e}tale}}
\def\etale{{\acute{e}tale}}
\def\QCoh{\mathit{QCoh}}
\def\Ker{\mathop{\mathrm{Ker}}}
\def\Im{\mathop{\mathrm{Im}}}
\def\Coker{\mathop{\mathrm{Coker}}}
\def\Coim{\mathop{\mathrm{Coim}}}

% Boxtimes
%
\DeclareMathSymbol{\boxtimes}{\mathbin}{AMSa}{"02}

%
% Macros for moduli stacks/spaces
%
\def\QCohstack{\mathcal{QC}\!\mathit{oh}}
\def\Cohstack{\mathcal{C}\!\mathit{oh}}
\def\Spacesstack{\mathcal{S}\!\mathit{paces}}
\def\Quotfunctor{\mathrm{Quot}}
\def\Hilbfunctor{\mathrm{Hilb}}
\def\Curvesstack{\mathcal{C}\!\mathit{urves}}
\def\Polarizedstack{\mathcal{P}\!\mathit{olarized}}
\def\Complexesstack{\mathcal{C}\!\mathit{omplexes}}
% \Pic is the operator that assigns to X its picard group, usage \Pic(X)
% \Picardstack_{X/B} denotes the Picard stack of X over B
% \Picardfunctor_{X/B} denotes the Picard functor of X over B
\def\Pic{\mathop{\mathrm{Pic}}\nolimits}
\def\Picardstack{\mathcal{P}\!\mathit{ic}}
\def\Picardfunctor{\mathrm{Pic}}
\def\Deformationcategory{\mathcal{D}\!\mathit{ef}}
\usepackage{cleveref}
\crefname{originalreference}{original reference}{original references}
\Crefname{originalreference}{Original reference}{Original references}
\makeatletter
\newlabel{dualizing-lemma-local-cohomology-adjoint}{{dualizing-lemma-local-cohomology-adjoint}{source}{}{sourcepointer-1}{}}
\newlabel{dualizing-lemma-local-cohomology-adjoint@cref}{{[originalreference][0][]dualizing-lemma-local-cohomology-adjoint}{[1][0][]source}{}{}{}}
\makeatother
\crefname{originalreference}{original reference}{original references}
\Crefname{originalreference}{Original reference}{Original references}
\makeatletter
\newlabel{dualizing-section-local-cohomology}{{dualizing-section-local-cohomology}{source}{}{sourcepointer-2}{}}
\newlabel{dualizing-section-local-cohomology@cref}{{[originalreference][0][]dualizing-section-local-cohomology}{[1][0][]source}{}{}{}}
\makeatother
\crefname{originalreference}{original reference}{original references}
\Crefname{originalreference}{Original reference}{Original references}
\makeatletter
\newlabel{more-algebra-lemma-local-cohomology-closed}{{more-algebra-lemma-local-cohomology-closed}{source}{}{sourcepointer-3}{}}
\newlabel{more-algebra-lemma-local-cohomology-closed@cref}{{[originalreference][0][]more-algebra-lemma-local-cohomology-closed}{[1][0][]source}{}{}{}}
\makeatother
\crefname{originallemma}{original lemma}{original lemmas}
\Crefname{originallemma}{Original lemma}{Original lemmas}
\makeatletter
\newlabel{lemma-local-cohomology-is-local-cohomology}{{lemma-local-cohomology-is-local-cohomology}{source}{}{sourcepointer-4}{}}
\newlabel{lemma-local-cohomology-is-local-cohomology@cref}{{[originallemma][0][]lemma-local-cohomology-is-local-cohomology}{[1][0][]source}{}{}{}}
\makeatother
\crefname{originallemma}{original lemma}{original lemmas}
\Crefname{originallemma}{Original lemma}{Original lemmas}
\makeatletter
\newlabel{lemma-cd}{{lemma-cd}{source}{}{sourcepointer-5}{}}
\newlabel{lemma-cd@cref}{{[originallemma][0][]lemma-cd}{[1][0][]source}{}{}{}}
\makeatother
\crefname{originalreference}{original reference}{original references}
\Crefname{originalreference}{Original reference}{Original references}
\makeatletter
\newlabel{coherent-lemma-vanishing-nr-affines}{{coherent-lemma-vanishing-nr-affines}{source}{}{sourcepointer-6}{}}
\newlabel{coherent-lemma-vanishing-nr-affines@cref}{{[originalreference][0][]coherent-lemma-vanishing-nr-affines}{[1][0][]source}{}{}{}}
\makeatother
\title{Three upper bounds for the cohomological dimension of a finitely generated ideal}
\author{The Stacks Project Authors}
\date{stacks-project:local-cohomology.tex:local-snapshot-20261012; GNU FDL 1.2 or later}
\begin{document}
\maketitle
\noindent\textit{Editorial note.} All delivered mathematical text is the author's own, copied from the pinned Stacks Project chapter source sources/stacks-local-cohomology/pinned-original-source.tex; no statement or proof step was rewritten, reordered or removed.\par\medskip
\noindent\textit{Editorial note.} Preamble provenance: the Stacks Project keeps one single shared preamble for every chapter; this card takes it directly from the staged sources/stacks-local-cohomology/preamble.tex as three verbatim spans (lines 8-12, 14-18 and 154-217), so no derived file was written and nothing was added or redefined. The wrapper class and the base packages are supplied by the assembler.\par\medskip
\noindent\textit{Editorial note.} Bibliography provenance: neither the counted statement nor its proof contains a citation (this lane rescans the whole span before writing the card), so the card carries no bibliography and no entry was reconstructed.\par\medskip
\noindent\textit{Editorial note.} Selection provenance: this card was opened under the file-based control contract (CONTROL.md), with the mechanical pre-card gate run on its statement and proof interval before the card was written.\par\medskip
\noindent\textit{Allowed standard prerequisites (editorial).}\begin{itemize}
\item The explicit description of the derived sections with support in the closed subset of a finitely generated ideal by the Cech type complex on the generators. Invoked at the start of the counted proof (source lines 705-706) as a cross-chapter pointer. Shows directly that the number of generators bounds the cohomological dimension.
\item The independence of the derived sections with support of the chosen finitely generated ideal defining the closed subset, together with the vanishing of the cohomology of a quasi-coherent sheaf on a space covered by a bounded number of affines. Both are invoked inside the counted proof (source lines 707-715 and 716-718) as cross-chapter pointers. The first transfers the bound from one set of generators to any other, the second converts an affine covering bound into a vanishing bound.
\end{itemize}\medskip
\ifdefined\hypertarget\hypertarget{sourcepointer-1}{}\fi
\noindent\textit{Source reference.} Original reference dualizing-lemma-local-cohomology-adjoint: sources/stacks-local-cohomology/pinned-original-source.tex, source lines 335--335. Navigation pointer only: the counted text refers to the Stacks item dualizing-lemma-local-cohomology-adjoint, which belongs to another chapter of the project and is not part of this card. The pointer carries the target's own label as its number because this chapter's staged snapshot does not contain the other chapter's tag table; its content is not reproduced here.\par\medskip
\ifdefined\hypertarget\hypertarget{sourcepointer-2}{}\fi
\noindent\textit{Source reference.} Original reference dualizing-section-local-cohomology: sources/stacks-local-cohomology/pinned-original-source.tex, source lines 241--241. Navigation pointer only: the counted text refers to the Stacks item dualizing-section-local-cohomology, which belongs to another chapter of the project and is not part of this card. The pointer carries the target's own label as its number because this chapter's staged snapshot does not contain the other chapter's tag table; its content is not reproduced here.\par\medskip
\ifdefined\hypertarget\hypertarget{sourcepointer-3}{}\fi
\noindent\textit{Source reference.} Original reference more-algebra-lemma-local-cohomology-closed: sources/stacks-local-cohomology/pinned-original-source.tex, source lines 714--714. Navigation pointer only: the counted text refers to the Stacks item more-algebra-lemma-local-cohomology-closed, which belongs to another chapter of the project and is not part of this card. The pointer carries the target's own label as its number because this chapter's staged snapshot does not contain the other chapter's tag table; its content is not reproduced here.\par\medskip
\ifdefined\hypertarget\hypertarget{sourcepointer-4}{}\fi
\noindent\textit{Source reference.} Original lemma lemma-local-cohomology-is-local-cohomology: sources/stacks-local-cohomology/pinned-original-source.tex, source lines 405--405. Navigation pointer only: the counted text refers to the Stacks item lemma-local-cohomology-is-local-cohomology, which belongs to another chapter of the project and is not part of this card. The pointer carries the target's own label as its number because this chapter's staged snapshot does not contain the other chapter's tag table; its content is not reproduced here.\par\medskip
\ifdefined\hypertarget\hypertarget{sourcepointer-5}{}\fi
\noindent\textit{Source reference.} Original lemma lemma-cd: sources/stacks-local-cohomology/pinned-original-source.tex, source lines 679--679. Navigation pointer only: the counted text refers to the Stacks item lemma-cd, which belongs to another chapter of the project and is not part of this card. The pointer carries the target's own label as its number because this chapter's staged snapshot does not contain the other chapter's tag table; its content is not reproduced here.\par\medskip
\ifdefined\hypertarget\hypertarget{sourcepointer-6}{}\fi
\noindent\textit{Source reference.} Original reference coherent-lemma-vanishing-nr-affines: sources/stacks-local-cohomology/pinned-original-source.tex, source lines 718--718. Navigation pointer only: the counted text refers to the Stacks item coherent-lemma-vanishing-nr-affines, which belongs to another chapter of the project and is not part of this card. The pointer carries the target's own label as its number because this chapter's staged snapshot does not contain the other chapter's tag table; its content is not reproduced here.\par\medskip
\begin{lemma}
\label{lemma-bound-cd}
Let $I \subset A$ be a finitely generated ideal of a ring $A$.
Then
\begin{enumerate}
\item $\text{cd}(A, I)$ is at most equal to the number of
generators of $I$,
\item $\text{cd}(A, I) \leq r$ if there exist $f_1, \ldots, f_r \in A$
such that $V(f_1, \ldots, f_r) = V(I)$,
\item $\text{cd}(A, I) \leq c$ if $\Spec(A) \setminus V(I)$
can be covered by $c$ affine opens.
\end{enumerate}
\end{lemma}

\begin{proof}
The explicit description for $R\Gamma_Y(-)$ given in
Dualizing Complexes, Lemma \ref{dualizing-lemma-local-cohomology-adjoint}
shows that (1) is true. We can deduce (2) from (1) using the
fact that $R\Gamma_Z$ depends only on the closed subset
$Z$ and not on the choice of the finitely generated ideal
$I \subset A$ with $V(I) = Z$. This follows either from the
construction of local cohomology in
Dualizing Complexes, Section \ref{dualizing-section-local-cohomology}
combined with
More on Algebra, Lemma \ref{more-algebra-lemma-local-cohomology-closed}
or it follows from Lemma \ref{lemma-local-cohomology-is-local-cohomology}.
To see (3) we use Lemma \ref{lemma-cd}
and the vanishing result of Cohomology of Schemes, Lemma
\ref{coherent-lemma-vanishing-nr-affines}.
\end{proof}

\end{document}
